\section{Abschließende Bewertung zur Eignung von modernen Dash\-boards}
% Abschließende Bewertung der Dashboard-Lösung zur Modernisierung eines Legacy-Systems
\label{ch:Diskussion:Abschließende_Bewertung}

% = Zusammenfassung aus Ergebnissem, Diskussion + Methodenkritik (wie Aussagekräftig die Ergebnisse sind)\\\\
Zusammenfassend lässt sich sagen, dass das Dash\-board eine geeignete Lösung zur Modernisierung des Legacy-Testpoolmonitoring-Prozesses darstellt. Die Ergebnisse der Usability-Auswertung stellen eine signifikante Verbesserung fest.

% Frage: Inwiefern beeinflusst die Nutzung eines modernen Dashboards das Monitoring der Ergebnisse des Legacy-Testpools?\\
Trotz der Verminderung der Aussagekraft des Usability-Tests bestätigen die Nutzer die Verbesserung anhand der Nutzerbefragung. Das Dash\-board hilft dabei, einen Überblick über alle Ergebnisse zu gewinnen. Da davon ausgegangen wird, dass mehr Nicht-Fehler als Fehler angezeigt werden, verringert sich der Aufwand. Für die Detailübersicht wird im Fehlerfall jeweils nur ein zusätzlicher  Mausklick benötigt.

Ein Problem stellt dar, dass das neue Dash\-board bisher nur ein Prototyp ist, um zu zeigen, dass die vorgeschlagene Modernisierungsmaßnahme gut begründet ist. Zur Vervollständigung werden weitere Ressourcen in Form von Personal und Zeit benötigt, die jetzt bereits knapp sind (s. \autoref{ch:SdF:Legacy}). Allerdings sind Modernisierungsprozesse immer aufwändig, wodurch dieses Argument im Speziellen gegen die vorgeschlagene Dash\-board"=Lösung hinfällig ist.

% Pro / Vorteile - Wie kann das Dashboard überzeugen?
Da das alte System nicht direkt abgelöst wird, müssen potenzielle Verbesserungen und Features in beiden Systemen umgesetzt werden. Um nicht selbst zur nächsten großen Wartungsaufgabe zu werden, ist der Parser zum Erhalt der Daten so geschrieben, dass das alte System unabhängig von den Inhalten ausgelesen wird (s. \autoref{Add:Dashboard:Parser_Code}). Sollten keine großen strukturellen Veränderungen am Altsystem erfolgen, was angesichts des jahrelangen Stillstands dessen unwahrscheinlich wirkt, benötigt das Dash\-board diesbezüglich wenig Wartungsaufwand.

Ein weiterer Vorteil der Lösung ist, dass Personen, die Jahrzehnte mit dem alten System gearbeitet haben, auch weiter mit dem alten System arbeiten können, da dieses vollständig unberührt bleibt (s. \autoref{ch:Dashboard:sec:Datenfluss}).

Da das Dash\-board nicht in das bestehende System eingreift, sondern nur darauf aufsitzt, ist es eine nicht invasive Lösung mit geringem Risiko in der Umsetzung. Es eignet sich als erster Schritt zur Verbesserung der Usability des Monitorings. Lange aufgeschobene Probleme und Altlasten, die sich in ihrer Komplexität potenzieren, konnten aufgedeckt werden, damit sie doch noch gelöst werden (s. \autoref{ch:Ergebnisse:Intuitivität}).

Des Weiteren kann es in mehreren Bereichen im Vergleich zum Altsystem mit neuen Funktionen überzeugen (s. \autoref{Add:ISO-Tabellen}, \autoref{ch:Dashboard:sec:Anforderungen_Auswertung}).
Es konnten die Ergebnisse von allen vorherigen Oberflächen auf wenige vereint werden. Dabei wurde die Struktur vereinfacht und selbsterklärender, indem die Ergebnisse in Fehlerkategorien zusammengefasst, Fachsprache simplifiziert und nicht benötigte Funktionen entfernt wurden. Zudem überzeugt das Dash\-board mit einer einheitlicheren, moderneren und intuitiveren Oberflächengestaltung, welche zusätzlich ergänzende Informationen liefert.

%Empfehlungen ableiten, was beim Übergang alt-neu wichtig ist -> hab ich glaub ich

% Wie gut eignen sich moderne Dashboards als benutzerfreundliche Modernisierung für das Monitoring der Ergebnisse des Legacy-Testpools?\\
Die Einführung des Dash\-boards ist demnach eine sehr sinnvolle Methode, um den Legacy-Testpoolmonitoring-Prozess effizienter und nutzerfreundlicher zu gestalten. %Trotz der positiven Rückmeldungen der Teilnehmer stellt diese Einführung des Dashboards jedoch nur einen Anfang des Modernisierungsprozesses dar.

%FF: Wie sinnvoll ist die Einführung eines modernen Dashboards im Hinblick auf das Ziel, den Legacy-Testpoolmonitoring-Prozess effizienter und nutzerfreundlicher zu gestalten?
